\input{hand-preamble.tex}
\newcommand{\gr}{\text{°}}
\newcommand{\ug}{\angle}
% дуга угла от руки
\newcommand{\harc}[4]{\draw[line width=0.6pt, decorate, decoration={random steps, segment length=5pt, amplitude=0.3pt}] ($(#1)+(#2:#4)$) arc (#2:#3:#4);}
\begin{document}\shapka
{\fontsize{20pt}{10mm}\selectfont Задание 1. Запись доказательства}\par
Теорема о сумме углов треугольника \framebox{Атанасян, 7--9, п. 30}\par\vspace{2mm}
\term{Категорическая форма.}\par
Сумма углов треугольника равна $180\gr$.\par
\term{Условная (импликативная) форма.}\par
Если $ABC$ --- треугольник, то $\ug A+\ug B+\ug C=180\gr$.\par\vspace{1mm}
\noindent\begin{minipage}[t]{0.5\textwidth}
\term{Дано:}\par
1. $\triangle ABC$\par
\term{Доказать:}\par
2. $\ug A+\ug B+\ug C=180\gr$\par
\nb{обозн.: $\ug 1=\ug A$, $\ug 2=\ug B$,}\par\nb{\hspace*{12mm}$\ug 3=\ug C$}
\end{minipage}\hfill
\begin{minipage}[t]{0.47\textwidth}\vspace{-6mm}
\begin{tikzpicture}[scale=1.05]
\coordinate (A) at (0,0); \coordinate (C) at (5.6,0.05); \coordinate (B) at (2.1,2.9);
\handseg{(A)}{(C)}\handseg{(A)}{(B)}\handseg{(B)}{(C)}
\handseg{(-0.6,2.88)}{(6.4,2.98)}
\node at (6.3,3.3) {$a$};
\node at (-0.3,-0.3) {$A$}; \node at (5.9,-0.25) {$C$}; \node at (2.05,3.35) {$B$};
\harc{A}{2}{54}{0.6} \node at (0.95,0.32) {\nb{1}};
\harc{B}{-126}{-55}{0.5} \node at (2.2,2.05) {\nb{2}};
\harc{C}{142}{178}{0.75} \node at (4.6,0.3) {\nb{3}};
\harc{B}{181}{233}{0.55} \node at (1.25,2.55) {\nb{4}};
\harc{B}{-52}{-1}{0.65} \node at (3.0,2.5) {\nb{5}};
\end{tikzpicture}
\end{minipage}\par\vspace{2mm}
\term{Доказательство:}\par
\def\st{\par\hangindent=7mm\hangafter=1\noindent}
\st 3. Через $B$ проведём прямую $a\parallel AC$ (доп. построение; через точку вне прямой можно провести прямую, параллельную данной).
\st 4. $\ug 4=\ug 1$ (накрест лежащие углы при $a\parallel AC$ и секущей $AB$; св-во параллельных прямых, 3).
\st 5. $\ug 5=\ug 3$ (накрест лежащие углы при $a\parallel AC$ и секущей $BC$; св-во параллельных прямых, 3).
\st 6. $\ug 4+\ug 2+\ug 5=180\gr$ (углы $4$, $2$, $5$ образуют развёрнутый угол с вершиной $B$, 3).
\st 7. $\ug 1+\ug 2+\ug 3=180\gr$ (6, 4, 5).
\st 8 (2). $\ug A+\ug B+\ug C=180\gr$ (7, обозн.). $\blacksquare$\par
\end{document}
